\section{Size the canonical header before storing it}
The canonical header has a fixed nine-byte contribution plus the decimal digit counts
of width and height. Each digit count is one or two under the exercise bounds. Thus
\[
  L=9+\operatorname{digits}(w)+\operatorname{digits}(h)+3wh.
\]
For the worked image, \(L=23\). At sixty-four by sixty-four it is 12301 bytes.
Validate dimensions first, then compute this bounded length and compare it with capacity
before any pixel read or header store. A short destination should not contain a new
magic number followed by an old raster.

Write the magic and LF, width, space, height, LF, then \texttt{255} and LF. For a
two-digit dimension, unsigned division by ten supplies the tens digit and remainder
supplies the units digit. Add ASCII zero to each numerical digit. A one-digit dimension
omits the zero quotient rather than writing a leading zero. Width nine and width ten
therefore exercise different header lengths even for the same height.

Once the header is complete, extract one input word's red, green and blue fields into
three consecutive bytes. The high byte does not enter the file. A byte-oriented encoder
can instead load the appropriate three byte positions from each little-endian word.
Both methods must preserve the sample codes and emit the same canonical bytes. The
completed-interface checks should accept either method.
